% PCC – Metodologia
\section{Metodologia}
A presente pesquisa caracteriza-se como uma pesquisa aplicada e experimental, com abordagem quantitativa para a avaliação de resultados. O trabalho foca no desenvolvimento de uma ferramenta de \textit{fuzzing} de caixa-preta (\textit{black-box}) orientada pela semântica de contratos OpenAPI.

\subsection{Arquitetura da Solução}
A ferramenta será construída de forma modular em Python, composta pelos seguintes componentes principais baseados na estrutura do projeto:
\begin{itemize}
    \item \textbf{Parser (\texttt{openapi\_parser.py}):} Responsável por ler arquivos YAML/JSON e extrair rotas, métodos e as restrições (\textit{constraints}) de dados.
    \item \textbf{Motor de Mutações (\texttt{engine.py}):} Implementa a lógica de geração de \textit{payloads}, aplicando mutações sistemáticas em tipos primitivos e estruturas complexas.
    \item \textbf{Orquestrador e Executor (\texttt{runner.py}):} Gerencia o plano de execução e realiza as chamadas assíncronas ao servidor-alvo.
    \item \textbf{Analisador (\texttt{analyzer.py}):} Atua como o oráculo do teste, comparando a resposta da API com as definições do contrato.
\end{itemize}

\subsection{Etapas de Desenvolvimento}
O estudo será conduzido seguindo as sete etapas detalhadas abaixo:

\begin{enumerate}
    \item \textbf{Revisão Teórica:} Estudo aprofundado sobre a especificação OpenAPI, fundamentos de \textit{fuzzing} e análise de ferramentas do estado da arte como Schemathesis e Dredd.
    \item \textbf{Modelagem do Contrato:} Desenvolvimento da lógica de extração de metadados do OpenAPI, convertendo definições estáticas em modelos de dados para teste.
    \item \textbf{Implementação dos Mutadores:} Criação de um catálogo de heurísticas para mutação de \textit{strings}, números, \textit{enums} e objetos, focando em violações de limites e tipos.
    \item \textbf{Desenvolvimento do Executor:} Implementação da \textit{engine} de requisições com suporte a concorrência e tratamento de erros de rede.
    \item \textbf{Módulo de Análise:} Desenvolvimento de heurísticas para identificar inconsistências (erros 5xx) e divergências estruturais entre a implementação e o contrato.
    \item \textbf{Avaliação Experimental:} Execução de testes em APIs de controle e cenários reais para medir a taxa de detecção de falhas e cobertura de \textit{endpoints}.
    \item \textbf{Discussão e Conclusão:} Análise crítica dos resultados obtidos, validação da hipótese de pesquisa e documentação de limitações da ferramenta.
\end{enumerate}
